\section{Truncation of complexes}
\label{section-truncations}

\noindent
Let $\mathcal{A}$ be an abelian category.
Let $A_\bullet$ be a chain complex. There
are several ways to {\it truncate} the complex $A_\bullet$.
\begin{enumerate}
\item The {\it ``stupid'' truncation $\sigma_{\leq n}$}
is the subcomplex $\sigma_{\leq n} A_\bullet$ defined
by the rule $(\sigma_{\leq n} A_\bullet)_i = 0$ if
$i > n$ and $(\sigma_{\leq n} A_\bullet)_i = A_i$ if
$i \leq n$. In a picture
$$
\xymatrix{
\sigma_{\leq n}A_\bullet \ar[d]  &
\ldots \ar[r] &
0 \ar[r] \ar[d] &
A_n \ar[r] \ar[d] &
A_{n - 1} \ar[r] \ar[d] &
\ldots \\
A_\bullet  &
\ldots \ar[r] &
A_{n + 1} \ar[r] &
A_n \ar[r] &
A_{n - 1} \ar[r] &
\ldots
}
$$
Note the property
$\sigma_{\leq n}A_\bullet / \sigma_{\leq n - 1}A_\bullet = A_n[-n]$.
\item The {\it ``stupid'' truncation $\sigma_{\geq n}$}
is the quotient complex $\sigma_{\geq n} A_\bullet$ defined
by the rule $(\sigma_{\geq n} A_\bullet)_i = A_i$ if
$i \geq n$ and $(\sigma_{\geq n} A_\bullet)_i = 0$ if
$i < n$. In a picture
$$
\xymatrix{
A_\bullet \ar[d]  &
\ldots \ar[r] &
A_{n + 1} \ar[r] \ar[d] &
A_n \ar[r] \ar[d] &
A_{n - 1} \ar[r] \ar[d] &
\ldots \\
\sigma_{\geq n}A_\bullet  &
\ldots \ar[r] &
A_{n + 1} \ar[r] &
A_n \ar[r] &
0 \ar[r] &
\ldots
}
$$
The map of complexes
$\sigma_{\geq n}A_\bullet \to \sigma_{\geq n + 1}A_\bullet$ is surjective
with kernel $A_n[-n]$.
\item The {\it canonical truncation} $\tau_{\geq n}A_\bullet$
is defined by the picture
$$
\xymatrix{
\tau_{\geq n}A_\bullet \ar[d]  &
\ldots \ar[r] &
A_{n + 1} \ar[r] \ar[d] &
\Ker(d_n) \ar[r] \ar[d] &
0 \ar[r] \ar[d] &
\ldots \\
A_\bullet  &
\ldots \ar[r] &
A_{n + 1} \ar[r] &
A_n \ar[r] &
A_{n - 1} \ar[r] &
\ldots
}
$$
Note that these complexes have the property that
$$
H_i(\tau_{\geq n}A_\bullet) =
\left\{
\begin{matrix}
H_i(A_\bullet) & \text{if} & i \geq n \\
0 & \text{if} & i < n
\end{matrix}
\right.
$$
\item The {\it canonical truncation} $\tau_{\leq n}A_\bullet$
is defined by the picture
$$
\xymatrix{
A_\bullet \ar[d]  &
\ldots \ar[r] &
A_{n + 1} \ar[r] \ar[d] &
A_n \ar[r] \ar[d] &
A_{n - 1} \ar[r] \ar[d] &
\ldots \\
\tau_{\leq n}A_\bullet  &
\ldots \ar[r] &
0 \ar[r] &
\Coker(d_{n + 1}) \ar[r] &
A_{n - 1} \ar[r] &
\ldots
}
$$
Note that these complexes have the property that
$$
H_i(\tau_{\leq n}A_\bullet) =
\left\{
\begin{matrix}
H_i(A_\bullet) & \text{if} & i \leq n \\
0 & \text{if} & i > n
\end{matrix}
\right.
$$
\end{enumerate}

\noindent
Let $\mathcal{A}$ be an abelian category.
Let $A^\bullet$ be a cochain complex. There
are four ways to truncate the complex $A^\bullet$.
\begin{enumerate}
\item The {\it ``stupid'' truncation $\sigma_{\geq n}$} is the subcomplex
$\sigma_{\geq n} A^\bullet$ defined by the rule
$(\sigma_{\geq n} A^\bullet)^i = 0$ if
$i < n$ and $(\sigma_{\geq n} A^\bullet)^i = A^i$ if
$i \geq n$. In a picture
$$
\xymatrix{
\sigma_{\geq n}A^\bullet \ar[d]  &
\ldots \ar[r] &
0 \ar[r] \ar[d] &
A^n \ar[r] \ar[d] &
A^{n + 1} \ar[r] \ar[d] &
\ldots \\
A^\bullet  &
\ldots \ar[r] &
A^{n - 1} \ar[r] &
A^n \ar[r] &
A^{n + 1} \ar[r] &
\ldots
}
$$
Note the property
$\sigma_{\geq n}A^\bullet / \sigma_{\geq n + 1}A^\bullet
= A^n[-n]$.
\item The {\it ``stupid'' truncation $\sigma_{\leq n}$}
is the quotient complex $\sigma_{\leq n} A^\bullet$ defined
by the rule $(\sigma_{\leq n} A^\bullet)^i = 0$ if
$i > n$ and $(\sigma_{\leq n} A^\bullet)^i = A^i$ if
$i \leq n$. In a picture
$$
\xymatrix{
A^\bullet \ar[d]  &
\ldots \ar[r] &
A^{n - 1} \ar[r] \ar[d] &
A^n \ar[r] \ar[d] &
A^{n + 1} \ar[r] \ar[d] &
\ldots \\
\sigma_{\leq n}A^\bullet &
\ldots \ar[r] &
A^{n - 1} \ar[r] &
A^n \ar[r] &
0 \ar[r] &
\ldots \\
}
$$
The map of complexes
$\sigma_{\leq n}A^\bullet \to \sigma_{\leq n - 1}A^\bullet$ is surjective
with kernel $A^n[-n]$.
\item The {\it canonical truncation} $\tau_{\leq n}A^\bullet$
is defined by the picture
$$
\xymatrix{
\tau_{\leq n}A^\bullet \ar[d]  &
\ldots \ar[r] &
A^{n - 1} \ar[r] \ar[d] &
\Ker(d^n) \ar[r] \ar[d] &
0 \ar[r] \ar[d] &
\ldots \\
A^\bullet  &
\ldots \ar[r] &
A^{n - 1} \ar[r] &
A^n \ar[r] &
A^{n + 1} \ar[r] &
\ldots
}
$$
Note that these complexes have the property that
$$
H^i(\tau_{\leq n}A^\bullet) =
\left\{
\begin{matrix}
H^i(A^\bullet) & \text{if} & i \leq n \\
0 & \text{if} & i > n
\end{matrix}
\right.
$$
\item The {\it canonical truncation} $\tau_{\geq n}A^\bullet$
is defined by the picture
$$
\xymatrix{
A^\bullet \ar[d] &
\ldots \ar[r] &
A^{n - 1} \ar[r] \ar[d] &
A^n \ar[r] \ar[d] &
A^{n + 1} \ar[r] \ar[d] &
\ldots \\
\tau_{\geq n}A^\bullet &
\ldots \ar[r] &
0 \ar[r] &
\Coker(d^{n - 1}) \ar[r] &
A^{n + 1} \ar[r] &
\ldots
}
$$
Note that these complexes have the property that
$$
H^i(\tau_{\geq n}A^\bullet) =
\left\{
\begin{matrix}
0 & \text{if} & i < n \\
H^i(A^\bullet) & \text{if} & i \geq n
\end{matrix}
\right.
$$
\end{enumerate}


















